% ===================================================================
%  పట్టిక సంఖ్య 16 — పెద్ద దుకాణం. బజారు. ఆటబొమ్మలు.
%
%  Traduction, faite en 2026, en telougou d'aujourd'hui, sur l'ido de
%  texto/io. Regles de la colonne : voir 10-pattika-01.tex. Dernier
%  tableau du livret : dix-huit alineas, cent trois renvois dont six
%  en lettres.
%
%  LA LETTRE (d) EST REMISE. Le fac-simile ido compose « d) » sans sa
%  parenthese ouvrante et renvoji.py ne peut donc pas la voir : c'est
%  l'un des trois ecarts declares du controle. Le francais la grave
%  correctement, et la colonne l'ecrit bien formee, comme le font
%  toutes les colonnes traduites depuis la premiere.
%
%  SEPT RETOURNEMENTS SUR VINGT PAIRES, sur cent trois renvois. Le
%  bazar enumere comme le marche du tableau 15 : la menagerie en
%  carton, les instruments de musique, les fleurs artificielles, les
%  animaux etranges sont des listes, et une liste ne s'inverse pas.
%  Treize des vingt paires signalees sont des membres d'enumeration
%  que la liste ne pouvait pas distinguer d'un rattachement — c'est
%  le prix de l'outil, et il est bas.
%
%  LA MENAGERIE EN CARTON EST LE MEILLEUR TEST DE VOCABULAIRE DU
%  LIVRET, parce qu'elle nomme vingt animaux d'un coup, et le
%  telougou en nomme la quasi-totalite sans emprunter : ఖడ్గమృగం le
%  rhinoceros, నిప్పుకోడి l'autruche, దుమ్ములగొండి la hyene,
%  నీటిగుర్రం l'hippopotame, ఒంటె le chameau, మొసలి le crocodile,
%  కొండచిలువ le boa, తిమింగలం la baleine, సొరచేప le requin, గబ్బిలం
%  la chauve-souris, తాబేలు la tortue, తోడేలు le loup. Trois seuls
%  manquent, et ce sont les trois qui ne vivent pas sous ces
%  latitudes : రెయిన్‌డీర్ le renne, వీసెల్ la belette, సీల్ le
%  phoque. La regle du marron chaud, une derniere fois, et sur un
%  echantillon assez large pour ne pas etre un hasard.
%
%  DEUX FELINS QU'IL FALLAIT SEPARER. Le livret donne « pantero »
%  (15) et « leopardo » (18) comme deux betes distinctes ; le
%  telougou courant appelle les deux చిరుతపులి. La colonne ecrit
%  నల్ల చిరుత pour la panthere — la panthere noire, qui est bien
%  l'animal des menageries de 1926 — et garde చిరుతపులి pour le
%  leopard. Deux renvois, deux mots : le lecteur qui suit les deux
%  colonnes doit trouver deux betes, pas deux fois la meme.
% ===================================================================

%%K t16-apar-1 apar
\VUsaut{4.0mm}
\VUtitre{15.0pt}{60}{పట్టిక సంఖ్య 16}
\VUsaut{2.0mm}
\VUpk{13.2pt}{40}{పెద్ద దుకాణం. --- బజారు.}
\VUpk{13.2pt}{40}{ఆటబొమ్మలు.}
\VUsaut{3.4mm}

%%K t16-01-1 p
1. --- కొత్త సంవత్సరం దగ్గరపడుతోంది. పెద్ద దుకాణాలు \VUgras{ఆటబొమ్మల}
\textsuperscript{(1)} ప్రదర్శనలూ కొత్త సంవత్సరపు కానుకలూ సిద్ధం చేశాయి.

%%K t16-01-2 p
ఆ అందమైన ప్రదర్శన చూసేందుకు బజారుకు తీసుకెళ్ళమని మా తాతయ్యను అడిగాను, ఆయన ఒప్పుకున్నారు.

%%K t16-02-1 p
2. --- మొదట అందమైన ఆయుధాల పలకల ముందు ఆగాం; వాటిలో \VUgras{తుపాకి},
\VUgras{సైనిక టోపీ}, \VUgras{వంపు కత్తి}, \VUgras{ఖడ్గపు నడికట్టు}, జత
\VUgras{భుజ పట్టీలు}, లేదా \VUgras{శిరస్త్రాణం}, \VUgras{కవచం} \textsuperscript{(3)},
\VUgras{పిస్తోలు} \textsuperscript{(4)} ఉన్నాయి. ఇవన్నీ \VUgras{వెలుగు ప్రొజెక్టర్ల}
\textsuperscript{(2)} కంటే నాకు చాలా ఎక్కువ ఆసక్తి కలిగించాయి.

%%K t16-tit-1 sub
\VUsaut{3.0mm}
\VUpk{10.2pt}{40}{అందమైన దృశ్యాలు.}
\VUsaut{2.6mm}

%%K t16-03-1 p
3. --- అదే విభాగంలో తన \VUgras{కాపలా గదిలో} \VUgras{కాపలా సైనికుడు}
\textsuperscript{(5)} ఉన్నాడు; అట్టతో చేసిన జంతుప్రదర్శనశాల ముందు కాపలా కాస్తున్నట్టు
కనిపించాడు. అక్కడ ఎన్ని జంతువులో : తన కొయ్యపై \VUgras{చిలుక} \textsuperscript{(6)},
ముక్కుపై కొమ్ముతో \VUgras{ఖడ్గమృగం} \textsuperscript{(7)}, \VUgras{రెయిన్‌డీర్}
\textsuperscript{(8)}, \VUgras{నిప్పుకోడి} \textsuperscript{(9)}, అక్రోటు కొరుకుతున్న
\VUgras{కోతి} \textsuperscript{(10)}, కోడిని ఎత్తుకుపోతున్న \VUgras{నక్క}
\textsuperscript{(11)}, భారీ దవడలు తెరిచిన \VUgras{మొసలి} \textsuperscript{(12)},
\VUgras{ఒంటె} \textsuperscript{(13)}, అందమైన \VUgras{సన్నని వేట కుక్క}
\textsuperscript{(14)}, \VUgras{నల్ల చిరుత} \textsuperscript{(15)}, ఆ తర్వాత నాకు
తెలియని రెండు జంతువులు — అవి \VUgras{వీసెల్} \textsuperscript{(16)},
\VUgras{దుమ్ములగొండి} \textsuperscript{(17)} అని చెప్పారు. అక్కడే \VUgras{చిరుతపులి}
\textsuperscript{(18)}, \VUgras{తోడేలు} \textsuperscript{(21)} కూడా ఉన్నాయి; చాలా
దగ్గరలో ముద్దైన \VUgras{ఎలుకలు} \textsuperscript{(19)}, రబ్బరుతో చేసిన అతి చిన్న
\VUgras{చిట్టెలుకలు} \textsuperscript{(20)} ఉన్నాయి. చివరికి \VUgras{నీటిగుర్రం}
\textsuperscript{(22)}, మూపురమున్న \VUgras{ఒంటి మూపురం ఒంటె} \textsuperscript{(23)} ఆ
సేకరణను ముగించాయి. పక్కనే \VUgras{సీసపు సైనిక బొమ్మలతో} \textsuperscript{(29)} నిండిన
అద్భుతమైన శిబిరం నా దృష్టిని లాగకపోతే నేను అక్కడ ఇంకా చాలా గంటలు ఉండేవాణ్ణి.

%%K t16-tit-2 sub
\VUsaut{4.0mm}
\VUpk{10.2pt}{40}{సైనికులు, యుద్ధం.}
\VUsaut{2.8mm}

%%K t16-04-1 p
4. --- మా తాతయ్య \VUgras{యుద్ధ వికలాంగుడు} \textsuperscript{(24)}; ఒక \VUgras{గాయం}
వల్ల ఆయన సైన్యం వదలాల్సి వచ్చింది, ఆ గాయమే ఆయనకు \VUgras{సైనిక పతకం} తెచ్చిపెట్టింది.
తాను పాల్గొన్న \VUgras{యుద్ధ యాత్రల} గురించి నాకు చెప్పడం ఆయనకు చాలా ఇష్టం. ఆ చిన్న
బొమ్మల సైన్యపు రకరకాల కదలికలు వివరిస్తూ ఆయన సంతోషపడ్డారు. బజారు చూడటానికి వచ్చిన నిజమైన
సైనికుల గుంపు — \VUgras{ఇంజనీరు దళ సైనికుడు} \textsuperscript{(25)},
\VUgras{బెరే టోపీ} పెట్టుకున్న \VUgras{పర్వత దళ సైనికుడు} \textsuperscript{(26)}, పదాతిదళ
\VUgras{సార్జెంటు మేజరు} \textsuperscript{(27)}, చేతిమీద బంగారు \VUgras{పట్టీతో}
\VUgras{ఫిరంగిదళ సార్జెంటు} \textsuperscript{(28)} — వివరణలు వినేందుకు మా పక్కన ఆగింది.

%%K t16-05-1 p
5. --- అంత మెరిసే శ్రోతలు దొరికినందుకు ఎంతో గర్వపడి, మా తాతయ్య ఒక పూర్తి యుద్ధ ప్రణాళికను
అప్పటికప్పుడు అల్లారు.

%%K t16-05-2 p
తన \VUgras{ప్రాకారాలతో}, \VUgras{కందకాలతో} ఉన్న కోట నగరాన్ని \VUgras{శత్రు సైన్యం}
ముట్టడించింది; సుదీర్ఘ \VUgras{ఫిరంగి వర్షం} తర్వాత అది \VUgras{దాడికి} సిద్ధమైంది.
కానీ ఆకలితో ఇబ్బంది పడ్డ \VUgras{ముట్టడిలో ఉన్నవాళ్ళు} \VUgras{ముట్టడించేవాళ్ళ}
\VUgras{శిబిరంపై} \VUgras{ఎదురుదాడికి} ప్రయత్నించారు. \VUgras{గుడారాల} చుట్టూ
పట్టువదలని \VUgras{పోరాటంతో} ఆ \VUgras{దాడిని} తిప్పికొట్టాక,
\VUgras{గెలిచిన} ముట్టడిదారులు \VUgras{ఓడిన} దాడిదారులను తరిమేందుకు తమ
\VUgras{అశ్వదళాలను} విడిచారు. వాళ్ళు \VUgras{వెనక్కి తగ్గారు},
\VUgras{యుద్ధ భూమిని} \VUgras{చనిపోయినవాళ్ళతో}, \VUgras{గాయపడినవాళ్ళతో} కప్పేస్తూ.
\VUgras{స్ట్రెచరు మోసేవాళ్ళు} ఆ చివరివాళ్ళను \VUgras{అంబులెన్సులోకి} ఎత్తుకెళ్ళారు —
\VUgras{శస్త్రవైద్యుడితో} చికిత్స చేయించేందుకు.

%%K t16-05-3 p
\VUgras{చతురస్రం} కట్టిన కొన్ని \VUgras{దళాల} వీరోచిత ప్రతిఘటన ఉన్నా \VUgras{ఓటమి}
పూర్తయింది; మిగిలిన సైన్యం \VUgras{పారిపోయింది}, చాలామంది \VUgras{బందీలను} వదిలేసి.
శత్రువు నగరాన్ని ఆక్రమించి తన \VUgras{విజయాన్ని} పూర్తి చేసుకోబోతున్నాడు. బహుశా అదే ఆ
యుద్ధ యాత్రకు ముగింపు కావచ్చు; \VUgras{కాల్పుల విరమణ} తర్వాత \VUgras{శాంతి ఒప్పందం}
కుదుర్చుకోవడం సాధ్యమవుతుంది.

%%K t16-tit-3 sub
\VUsaut{4.4mm}
\VUpk{10.2pt}{40}{కొత్త సంవత్సరపు కానుకలు.}
\VUsaut{2.8mm}

%%K t16-06-1 p
6. --- ఆ ముసలి సైనికుడి రమణీయమైన కథ వింటూ నవ్విన కుర్ర సైనికులు ఆయనను ఇంకొన్ని వివరాలు
అడిగారు. నా మటుకు నేను, అందమైన \VUgras{అడ్డ విల్లు} \textsuperscript{(32)},
\VUgras{విల్లు} \textsuperscript{(33)} చూసేందుకు దుకాణం చుట్టూ తిరిగాను — దాని
\VUgras{బాణంతో} \textsuperscript{(31)} సహా. అక్కడ ఇంకో జంతువుల సేకరణ కనిపించింది : రెండు
\VUgras{పాడే పక్షులు} \textsuperscript{(34)} — యంత్రంతో నడిచే \VUgras{నైటింగేల్},
గొంతెత్తి పాడే \VUgras{వార్బ్లరు} —, ఇంకా వరుసగా విచిత్రమైన జంతువులు :
\VUgras{గబ్బిలం} \textsuperscript{(35)}, \VUgras{కొండచిలువ} \textsuperscript{(36)},
\VUgras{తాబేలు} \textsuperscript{(37)}, \VUgras{సీల్} \textsuperscript{(38)},
\VUgras{తిమింగలం} \textsuperscript{(39)}, గుడ్డతో నింపిన \VUgras{సొరచేప}
\textsuperscript{(40)}.

%%K t16-07-1 p
7. --- వాటిని చూస్తూ నేను ఎక్కువసేపు ఆగలేదు — ఎందుకంటే నేను కలగన్నది కనిపించింది :
అద్భుతమైన \VUgras{బొమ్మలాట రంగస్థలం} \textsuperscript{(41)}, చక్కటి \VUgras{అలంకరణలతో},
మనసు దోచే ఎర్రని \VUgras{తెరతో}. \VUgras{రంగంపై} ఇద్దరు నటులు చాలా ఆసక్తికరమైన
\VUgras{నాటకంలో} \VUgras{పాత్ర} పోషిస్తున్నట్టు కనిపించారు — ఎందుకంటే అట్టతో చేసిన
చిన్న \VUgras{ప్రేక్షకులు}, గదుల్లో కూర్చున్నవాళ్ళూ, \VUgras{కుర్చీలపైన},
\VUgras{వాద్యబృంద నాయకుడి} వెనుక \VUgras{నేలభాగంలో} ఉన్నవాళ్ళూ ఉత్సాహంగా
చప్పట్లు కొడుతున్నారు.

%%K t16-07-2 p
దురదృష్టవశాత్తూ ఆ రంగస్థలం చాలా ఖరీదు.

%%K t16-08-1 p
8. --- పైగా ఆటబొమ్మలు ఎంచుకోవడం కష్టమే — అద్దాల అరల్లో అన్ని రకాల బొమ్మలూ కుప్పలుగా
ఉన్నాయి : \VUgras{ఆర్లెకిన్} \textsuperscript{(42)}, \VUgras{పులిచినెల్లో}
\textsuperscript{(43)}, \VUgras{పియెరో} \textsuperscript{(44)}, రెక్కలు టపటపలాడించే
\VUgras{గుడ్లగూబలు} \textsuperscript{(45)}, ఈలవేసే \VUgras{నల్ల పిట్టలు}
\textsuperscript{(46)}, మొరిగే \VUgras{పూడిల్ కుక్కలు}, \VUgras{గ్రామఫోను}
\textsuperscript{(47)}, \VUgras{ఓపిక ఆటలు}. వాటిలో ఒకటి నన్ను బాగా అలరించింది : అది

%%K t16-08-1 p suite
\VUgras{నౌకా యుద్ధాన్ని} \textsuperscript{(48)} చూపిస్తుంది — అందులో ఒక \VUgras{ఓడపై}
\VUgras{సముద్రపు దొంగలు} దాడి చేస్తారు, \VUgras{కొబ్బరి చెట్లు} నాటిన నేలకు దగ్గరగా.

%%K t16-noto-1 noto
\VUnotes{143.2mm}{%
\textit{(*) పట్టిక సంఖ్య 6 చూడండి.}%
}

%%K t16-09-1 p
9. --- ఆ అద్భుతాలన్నీ నేను చాలా తీరిగ్గా చూడగలిగాను — దుకాణంలో కొద్దిమందే ఉన్నారు,
అక్కడక్కడ కొందరు తిరుగుతున్నవాళ్ళు, \VUgras{డ్రాగన్ అశ్వసైనికుడు} \textsuperscript{(49)},
ఇంకా \VUgras{వసూలుదారు} \textsuperscript{(50)} — అతను ప్రధాన \VUgras{కాషు}
\textsuperscript{(51)} దగ్గర \VUgras{హుండీ} చూపిస్తున్నాడు.

%%K t16-10-1 p
10. --- \VUgras{కాషియరు} వెనుక బల్లలపై పెట్టిన పుస్తకాలు దేనికని మా తాతయ్యను అడిగాను;
అవి \VUgras{లెక్కల పుస్తకాలు} అని ఆయన చెప్పారు.

%%K t16-tit-4 sub
\VUsaut{3.6mm}
\VUpk{10.2pt}{40}{దుకాణం చుట్టూ నోరు తెరుచుకుని.}
\VUsaut{2.8mm}

%%K t16-11-1 p
11. --- ఆటబొమ్మల విభాగం వదిలి, జీను సామాను విభాగానికి వెళ్ళాం — \VUgras{జీనులు}
\textsuperscript{(53)}, \VUgras{రికాబులు} \textsuperscript{(54)}, \VUgras{కళ్ళేలు}
\textsuperscript{(55)}, \VUgras{కళ్ళెపు కడ్డీలు} \textsuperscript{(56)},
\VUgras{కొరడాలు} ఒకసారి చూసేందుకు. \VUgras{విభాగపు పెద్ద} \textsuperscript{(57)} మా
తాతయ్యకు కొన్ని వివరాలు చెబుతుండగా, నేను \VUgras{పింగాణీ},
\VUgras{స్ఫటిక పాత్రల} \textsuperscript{(58)} ప్రదర్శన చూడటానికి వెళ్ళాను — అక్కడ
అందమైన \VUgras{సలాడు గిన్నెలూ}, సున్నితమైన \VUgras{టీ కెటిల్లూ} \textsuperscript{(59)}
ఉన్నాయి.

%%K t16-12-1 p
12. --- మొదటి అంతస్తుకు వెళ్ళేందుకు \VUgras{కార్సెట్ల} \textsuperscript{(60)} అద్దాల అర
వెనుకనుంచి వెళ్ళాం — మేజోళ్ళ విభాగం పక్కన; అక్కడ చాలా అందమైన \VUgras{ప్యాంట్లు}
\textsuperscript{(63)} ప్రదర్శించారు. దగ్గరలో \VUgras{అమ్మకందారిణి} \textsuperscript{(61)}
\VUgras{కొనుగోలుదారిణితో} \textsuperscript{(62)} అందమైన పట్టు \VUgras{వస్త్రాలను}
\textsuperscript{(64)} ఎంపిక చేయిస్తోంది.

%%K t16-12-2 p
మెట్లు ఎక్కుతుండగా, దుకాణాన్ని జపాను \VUgras{లాంతర్లతో} \textsuperscript{(65)},
\VUgras{గొడుగులతో} \textsuperscript{(66)}, భారీ \VUgras{తాటి చెట్లతో} \textsuperscript{(70)}
అలంకరించారని గమనించాను.

%%K t16-13-1 p
13. --- మొదటి అంతస్తులో చూడదగ్గది ఎక్కువ లేదు; ఫర్నిచరు (*), \VUgras{ఇనుప పెట్టెలు}
\textsuperscript{(67)}, \VUgras{సొరుగుల బీరువాలు} \textsuperscript{(68)},
\VUgras{పిల్లల బండ్లు} \textsuperscript{(69)} మాత్రమే. కానీ దుకాణపు మొత్తం దృశ్యం మాత్రం
చాలా \VUgras{సరదాగా} ఉంది.

%%K t16-14-1 p
14. --- మా కాళ్ళ కింద సంగీత వాద్యాలు కనిపించాయి : \VUgras{హార్పు వీణలు}
\textsuperscript{(71)}, \VUgras{గిటార్లు} \textsuperscript{(72)},
\VUgras{మాండొలిన్లు} \textsuperscript{(73)}, \VUgras{పిస్టను బూరలు}
\textsuperscript{(74)}, \VUgras{క్లారినెట్లు} \textsuperscript{(75)}, వాటి
\VUgras{కమానుతో} \textsuperscript{(76)} వయొలిన్లు, చివరికి \VUgras{పెద్ద డోలు}
\textsuperscript{(77)} కూడా. కొంచెం అవతల, కాగితాల విభాగం దగ్గర, \VUgras{విద్యార్థి}
\textsuperscript{(78)} ఒకడు \VUgras{గుమాస్తాతో} \textsuperscript{(79)} నోటుపుస్తకాల డబ్బా
బేరమాడుతున్నాడు. వాళ్ళ పక్కన అందమైన \VUgras{అక్షరమాల పుస్తకం} \textsuperscript{(80)}
నాకు కనిపించింది.

%%K t16-14-2 p
దుకాణం మధ్యలో ఎత్తైన \VUgras{క్రిస్మస్ చెట్టు} \textsuperscript{(81)} నిలబెట్టారు —
కొవ్వొత్తులు, చిన్న వస్తువులు, అన్ని రకాల ఆటబొమ్మలతో నిండుగా : \VUgras{కర్ర గుర్రాలు}
\textsuperscript{(82)}, \VUgras{బొమ్మలు} \textsuperscript{(83)} ఇత్యాదులు.

%%K t16-14-3 p
\VUgras{సువాసన ద్రవ్యాల విభాగం} \textsuperscript{(84)} తన \VUgras{పంటి పొడులతో},
రకరకాల \VUgras{సువాసనలతో} చాలా మంచి వాసన వెదజల్లింది, అది మా దాకా వచ్చింది.

%%K t16-tit-5 sub
\VUsaut{3.4mm}
\VUpk{10.2pt}{40}{మేము కొంత కొంటాం.}
\VUsaut{2.8mm}

%%K t16-15-1 p
15. --- మా తాతయ్యకు సమయం ఎక్కువ అవుతోందని అనిపించడంతో, పెద్ద మెట్ల మీదుగా కిందికి దిగాం.
\VUgras{ఖగోళ శాస్త్ర పటం} \textsuperscript{(85)} ముందు ఆయన కాసేపు ఆగారు — అది
\VUgras{తోకచుక్కలు} \textsuperscript{(a)}, \VUgras{చంద్ర, సూర్య గ్రహణాలు}
\textsuperscript{(b)}, \VUgras{నాలుగు దిక్కులతో} \textsuperscript{(c)}
\VUgras{(ఉత్తరం, దక్షిణం, తూర్పు, పడమర)} దిక్చక్రం, రెండు \VUgras{అర్ధగోళాల}
\textsuperscript{(d)} పటం, \VUgras{గ్రహాలు} \textsuperscript{(e)},
\VUgras{సౌర వ్యవస్థ} \textsuperscript{(f)} చూపిస్తుందని ఆయన నాకు చెప్పారు. వీటిలో నాకు
ఏమీ అర్థం కాలేదు; పక్కనుంచి వెళ్తున్న \VUgras{సరుకు అందించే కుర్రాడి}
\textsuperscript{(86)} జరీ దుస్తులూ, నా పక్కనే ఉన్న అందమైన \VUgras{విసనకర్రలూ}
\textsuperscript{(87)} నాకు చాలా ఎక్కువ ఆసక్తి కలిగించాయి — ఆ విసనకర్రలు
\VUgras{డబ్బు సంచుల} \textsuperscript{(88)}, \VUgras{జేబు పర్సుల} \textsuperscript{(89)}
దగ్గర ఉన్నాయి.

%%K t16-16-1 p
16. --- ప్రదర్శన బల్లపై \VUgras{ఈకలు} \textsuperscript{(90)},
\VUgras{నడికట్టు బకిల్లు} \textsuperscript{(91)}, బుట్టల్లో ముచ్చటగా అమర్చిన అందమైన
\VUgras{కృత్రిమ పూలు} \textsuperscript{(94)} కూడా నేను మెచ్చుకున్నాను.
\VUgras{వైలెట్లు}, \VUgras{వాల్‌ఫ్లవర్లు}, \VUgras{హైసింత్‌లు}
\textsuperscript{(95)}, \VUgras{జెరేనియంలు} \textsuperscript{(96)},
\VUgras{లిల్లీ పూలు} \textsuperscript{(97)}, \VUgras{నాస్టర్షియంలు}
\textsuperscript{(98)} చాలా బాగా నకలు చేశారు.

%%K t16-tit-6 sub
\VUsaut{4.4mm}
\VUpk{10.2pt}{40}{మా తాతయ్య కానుక.}
\VUsaut{2.8mm}

%%K t16-17-1 p
17. --- \VUgras{పంటి బ్రష్షుల్లో} \textsuperscript{(92)} ఒకటి కొనిపెట్టమని మా తాతయ్యను
అడిగాను — అవి \VUgras{జుట్టు పిన్నుల} \textsuperscript{(93)} పక్కన ఒక పెట్టెలో ఉన్నాయి.
ఆయన ఒప్పుకున్నారు, నేనే స్వయంగా కాషు దగ్గరికి వెళ్ళి నా కొనుగోలుకు డబ్బు కట్టాను; అక్కడే
పెద్ద \VUgras{సరుకు రవాణా} \textsuperscript{(100)} సిద్ధం చేస్తున్నారు — ఒక
\VUgras{ఖాతాదారు} \textsuperscript{(99)} కోసం.

%%K t16-18-1 p
18. --- అకస్మాత్తుగా, కళాత్మక \VUgras{కంచు విగ్రహాల} \textsuperscript{(101)} దగ్గర ఆగిన
జనంలో చిన్న గందరగోళం రేగింది. ఎవరినో దోచుకోబోతుండగా \VUgras{దొంగను}
\textsuperscript{(102)} \VUgras{తనిఖీ అధికారి} \textsuperscript{(103)} ఇప్పుడే
రెడ్‌హ్యాండెడ్‌గా పట్టుకున్నాడు. అతన్ని \VUgras{అరెస్టు} చేసేందుకు పోలీసును
పిలిపించారు; మేము బయటికి వస్తున్న క్షణంలోనే ఆ నేరస్తుడిని చెరసాలకు తీసుకెళ్తున్నారు.

\VUsaut{6.0mm}
\VUfilet{22mm}
